\documentclass[11pt,a4paper]{article}

% ---------------------------------------------------------
% PREAMBLE (stable, warning-free, Overleaf-safe)
% ---------------------------------------------------------
\usepackage[utf8]{inputenc}
\usepackage[T1]{fontenc}
\usepackage{lmodern}
\usepackage{amsmath, amssymb, amsfonts}
\usepackage{bm}
\usepackage{geometry}
\usepackage{graphicx}
\usepackage{hyperref}
\usepackage{cite}
\usepackage{authblk}
\usepackage{physics}
\usepackage{mathtools}

\geometry{margin=2.4cm}

\title{\textbf{Companion LXXXIX: Hierarchical Bayesian RDCC Inference with Multi-Probe Priors and Transport-Regularized Posteriors}}
\author{Michael Lehmann \\ \texttt{mi.lehmann@gmx.de}}
\date{April 2026}

\begin{document}
\maketitle

\begin{abstract}
The Relaxation-Driven Cyclic Cosmology (RDCC) introduces the relaxation parameter 
$\alpha$, which modifies the scale-dependent gravitational potential and induces 
observable signatures in weak-lensing convergence fields. Previous Companions 
developed transport-based topological statistics, multi-probe likelihoods, and 
emulator-accelerated inference. In this work, we construct a hierarchical Bayesian 
framework for RDCC inference that incorporates multi-probe priors, 
transport-regularized posteriors, and emulator uncertainty. The hierarchical 
structure enables principled combination of topological, morphological, and 
two-point observables, while transport regularization ensures geometric 
consistency across probes. This Companion provides the final statistical layer 
required for robust, scalable, and fully Bayesian RDCC parameter estimation.
\end{abstract}
% ---------------------------------------------------------
\section{Introduction}

Weak-lensing surveys provide multiple complementary observables, including:

\begin{itemize}
    \item topological transport statistics,
    \item Minkowski functionals,
    \item shear power spectra,
    \item neural joint likelihoods (Companion LXXXVI),
    \item emulator-accelerated predictions (Companion LXXXVIII).
\end{itemize}

A hierarchical Bayesian framework is required to:

\begin{itemize}
    \item combine these probes coherently,
    \item propagate emulator uncertainty,
    \item incorporate survey systematics,
    \item regularize posteriors using geometric transport constraints.
\end{itemize}

This Companion constructs such a framework.
% ---------------------------------------------------------
\section{Hierarchical Model Structure}

We define the hierarchical model:



\[
    \alpha \sim p(\alpha),
\]





\[
    \theta_{\mathrm{emu}} \sim p(\theta_{\mathrm{emu}}),
\]





\[
    \mathbf{s} \sim p(\mathbf{s} \mid \alpha, \theta_{\mathrm{emu}}),
\]



where:

\begin{itemize}
    \item $\alpha$ is the RDCC relaxation parameter,
    \item $\theta_{\mathrm{emu}}$ are emulator parameters,
    \item $\mathbf{s}$ is the multi-probe summary vector.
\end{itemize}

The posterior is:



\[
    p(\alpha, \theta_{\mathrm{emu}} \mid \mathbf{s}_{\mathrm{obs}})
    \propto
    p(\mathbf{s}_{\mathrm{obs}} \mid \alpha, \theta_{\mathrm{emu}})
    p(\theta_{\mathrm{emu}})
    p(\alpha).
\]


% ---------------------------------------------------------
\section{Multi-Probe Priors}

We define a prior that incorporates:

\begin{itemize}
    \item theoretical RDCC constraints,
    \item emulator smoothness priors,
    \item cross-probe consistency priors.
\end{itemize}

The multi-probe prior is:



\[
    p(\alpha)
    \propto
    \exp\!\left[
        -\lambda_{\mathrm{th}}
        \left(
            \frac{\partial^{2} \mathcal{T}_{\ell}}{\partial \alpha^{2}}
        \right)^{2}
    \right],
\]



which penalizes unphysical curvature in transport statistics.
% ---------------------------------------------------------
\section{Transport-Regularized Posteriors}

Transport regularization enforces geometric consistency across probes:



\[
    \mathcal{R}_{\mathrm{trans}}
    =
    \sum_{\ell}
    W_{2}^{2}
    \left(
        S_{\theta}(\alpha),
        S_{\theta}(\alpha + \delta)
    \right),
\]



where $S_{\theta}$ is the neural transport surrogate.

The posterior becomes:



\[
    p(\alpha \mid \mathbf{s})
    \propto
    \mathcal{L}(\mathbf{s} \mid \alpha)
    \exp\!\left[
        -\lambda_{\mathrm{trans}}
        \mathcal{R}_{\mathrm{trans}}
    \right].
\]


% ---------------------------------------------------------
\section{Inference and Sampling}

The hierarchical posterior is sampled using:

\begin{itemize}
    \item Hamiltonian Monte Carlo,
    \item No-U-Turn Sampler (NUTS),
    \item Variational inference,
    \item Neural posterior estimation.
\end{itemize}

Gradients are computed via automatic differentiation through:

\begin{itemize}
    \item the neural joint likelihood (Companion LXXXVI),
    \item emulator flows (Companion LXXXVIII),
    \item transport regularization terms.
\end{itemize}
% ---------------------------------------------------------
\section{Conclusion}

We introduced a hierarchical Bayesian framework for RDCC inference that combines 
multi-probe priors, emulator uncertainty, and transport-regularized posteriors. 
This framework unifies topology, morphology, and two-point statistics into a 
single coherent Bayesian model and provides the final statistical layer required 
for robust RDCC parameter estimation in Euclid-like surveys.

\bibliographystyle{apsrev4-2}
\nocite{*}
\bibliography{references}


\end{document}
